\textbf{COLUNI 2019 -- Questão 05}\par\medskip Leia os fragmentos do poema a seguir:

\par

\textbf{O Fotógrafo}

\par

\begin{verse}
Difícil fotografar o silêncio.\\
Entretanto tentei. Eu conto:\\
[...]\\
Ia o Silêncio pela rua carregando um bêbado.\\
Preparei minha máquina.\\
O silêncio era um carregador?\\
Estava carregando o bêbado.\\
Fotografei esse carregador.\\
[...]\\
Tinha um perfume de jasmim no beiral de um sobrado.\\
Fotografei o perfume.\\
Vi uma lesma pregada na existência mais do que na pedra.\\
Fotografei a existência dela.\\
[...]\\
A foto saiu legal.
\end{verse}
(BARROS, Manoel de. O fotógrafo. In: \_\_\_\_\_. \textbf{Meu quintal é maior do que o mundo}. Rio de Janeiro: Alfaguara, 2015, p. 115.)

\par

Nesses versos identificamos um recurso linguístico muito utilizado na poesia de Manoel de Barros que consiste na:\begin{enumerate}
\item utilização, predominantemente, do relato de experiências sinestésicas, recriando um mundo de sonho.
\item reconstrução com exatidão de imagens e ações do passado, visando facilitar a compreensão do leitor.
\item apresentação de profusão sensorial, misturando-se realidade e imaginação no tempo e no espaço.
\item utilização da escrita memorialística da infância, demonstrando erudição e fluência vocabulares.\end{enumerate}
